%% Beispiel-Präsentation mit LaTeX Beamer im KIT-Design
%% entsprechend den Gestaltungsrichtlinien vom Februar 2025
%%
%% Siehe https://sdq.kastel.kit.edu/wiki/Dokumentvorlagen

%% Beispiel-Präsentation
\documentclass[en,navbaroff]{afvbeamer}
% kitgrid,de,navbaroff,smallfoot,helvet

\usepackage{hyphenat}
\usepackage{svg}
\usepackage{caption}
\usepackage{comment}
\usepackage{cancel}
\usepackage{lipsum}
\usepackage{tikz}
\usepackage{pgfplots}
\pgfplotsset{compat=1.16}
\usetikzlibrary {mindmap}
\usepackage{standalone}
\usepackage{esvect}
\usepackage{algorithm}
\usepackage{algpseudocode}

\usetikzlibrary{intersections}

%% Gruppenlogo, muss im Verzeichnis logos/ liegen
%% falls kein Gruppenlogo gewünscht, bitte \grouplogo{} aufrufen
\grouplogo{}

%% Gruppenname und Breite (Standard: 89 mm)
\groupname{}
%\groupnamewidth{89mm}

% Beginn der Präsentation

\title[Formalizing Thirty-Three Miniatures]{Formalizing Thirty-Three Miniatures}
\subtitle{Implementing Applications of Linear Algebra in Isabelle/HOL}
\author[David, Alicia]{David Kiefer \and Alicia Appelhagen}

\date[25.02.2026]{25.02.2026}

% Literatur 

\usepackage[citestyle=authoryear,bibstyle=numeric,hyperref,natbib,backend=biber]{biblatex}
\addbibresource{presentation.bib}
\bibhang1em

\newif\ifsol
\solfalse %\soltrue to show solutions, \solfalse to hide them

\newcommand{\negimplies}{%
	\mathrel{{\ooalign{\hidewidth$\neg\phantom{=}$\hidewidth\cr$\implies$}}}}
	
\tikzset{
	cross/.pic = {
		\draw[rotate = 45] (-#1,0) -- (#1,0);
		\draw[rotate = 45] (0,-#1) -- (0, #1);
	}
}

\begin{document}
	
	\begin{frame}[title white horizontal]%'kitlogo=white' for a darker image
		\titlepage
	\end{frame}
	
	\section{Motivation and Formalization Content}
	\begin{frame}{Formalization Goals}
		\begin{minipage}[h]{0.33\textwidth}
			\begin{orangeblock}{Formalizing Mathematics}
				\begin{itemize}
					\item Formalize proof strategies
					\item Focus on linear algebra methods
					\item Build on \cite{TODO book}
					\item Benefits of proof assistants
				\end{itemize}
			\end{orangeblock}
		\end{minipage}
		\begin{minipage}[h]{0.33\textwidth}
			\begin{blueblock}{Exploring Libraries}
				\begin{itemize}
					\item Build on existing formalizations
					\item Archive of Formal Proofs + Isabelle/HOL
					\item Extend linear algebra libraries
				\end{itemize}
			\end{blueblock}
		\end{minipage}
		\begin{minipage}[h]{0.33\textwidth}
			\begin{yellowblock}{Applying Isabelle/HOL}
				\begin{itemize}
					\item Understand Isabelle features
					\item Write reusable proofs
				\end{itemize}
			\end{yellowblock}
		\end{minipage}
	\end{frame}
	
	\begin{frame}{Formalizing Mathematics: Thirty-Three Miniatures}
		\framesubtitle{Miniature 2}
	\end{frame}
	
	\begin{frame}{Formalizing Mathematics: Thirty-Three Miniatures}
		\framesubtitle{Miniature 3}
	\end{frame}
	
	\begin{frame}{Formalizing Mathematics: Thirty-Three Miniatures}
		\framesubtitle{Thirty-Three Miniatures: Miniature 5}
	\end{frame}
	
	\begin{frame}{Formalizing Mathematics: Learnings}
		\begin{orangeblock}{Theorem Provers Encourage Abstraction}
			Abstraction makes proofs less verbose (less options for try, searching for relevant lemmas by hand is easier)
		\end{orangeblock}
		
		\begin{orangeblock}{Concrete Examples Can Help ($\#$instantiateYourLocales)}
			Concrete application of definitions avoids inconsistencies and hopeless proof attempts
			\begin{itemize}
				\item First oddtown definition did not require sets to be different in the oddtown predicate -> empty set, easy to prove?
				\item Hamming Code definition issues: universal quantification over code words instead of all words -> proof broke?
			\end{itemize}
		\end{orangeblock}
	\end{frame}

	
	\section{Exploring Libraries}
	
	\begin{frame}{Algebra Libraries in Isabelle/HOL}
		\begin{itemize}
			\item Algebraic structures as type classes (carrier = type universe) vs. locales (carrier explicitly given) [mention typedef, Rep and Abs]
			\item using ring-scheme for objects that do not have a multiplication
		\end{itemize}
	\end{frame}
	
	\begin{frame}{Algebra Libraries in the Archive of Formal Proofs}
		\begin{itemize}
			\item Algebraic structures as type classes (carrier = type universe) vs. locales (carrier explicitly given) [mention typedef, Rep and Abs]
			\item using ring-scheme for objects that do not have a multiplication
		\end{itemize}
	\end{frame}

	\section{Applying Isabelle/HOL}

	\begin{frame}{Useful Isabelle/HOL Features}
		\begin{itemize}
			\item Commands: proof-cases, try
			\item Types: records, typedef
			\item Update syntax: function updates (see Nitpick), record updates (subspace.vs)
			\item Locales (\enquote{locale1 + ...} leads to automatic instantiations of locale1-parameters and directly yields the assumptions of locale1)
			\item Proof structure: blocks (idk, the $\{...\}$ thingies) (according to some people this makes proofs less understandable so maybe not really that useful)
			\item <subgoal-proof> tactic applies tactic to the remaining subgoals like the \enquote*{proof(...)} command in the beginning
			\item meme: isabelle logo
		\end{itemize}
	\end{frame}
	
	\subsection{Eisbach}
	\begin{frame}{Eisbach}
		content... (TODO (cope))
	\end{frame}
	
	\subsection{Weird Stuff}
	\begin{frame}{Weird Stuff}
		\begin{itemize}
			\item leaving out \enquote{and}s in assumes/shows clauses changes the behaviour (leaving out the \enquote{shows} before multiple goals of a lemma leads to \enquote{thm} only referencing the first goal instead of all of them)
			\item lemmas that are unexpectedly not found by try0 (e.g. $n > 0 \Rightarrow \exists m. n = Suc \; m$)
			\item no error on forgotten \enquote{end} after theory/context
			\item locale instantiations with fixed parameters with no assumptions curried?
			\item \enquote{?} behind added locale \enquote*{name} makes \enquote*{name.def} optional (approved)
			\item list comprehension $[1::nat..3::nat]$ becomes $int \; list$ but $[1::nat..<3::nat]$ gives $nat \; list$
			\item Non-aligning AFP and Isabelle version
			\item termination proofs for non-recursive functions (defined by \enquote{function}) optional but if not provided, simps rules are whack
			\item termination using \enquote*{termination} by blast????
		\end{itemize}
	\end{frame}
	
\end{document}
